\documentclass[12pt,a4paper]{article}

\usepackage[margin=1in]{geometry}
\usepackage{setspace}
\usepackage[ngerman]{babel}
\usepackage{changepage}
\usepackage{newtxtext}
\usepackage[none]{hyphenat}
\usepackage{enumitem}
\usepackage[colorlinks=true,linkcolor=blue,urlcolor=blue,citecolor=blue]{hyperref}
\newcommand{\ulhref}[2]{\href{#1}{\uline{#2}}}


\pagestyle{empty}

\begin{document}
%
\begin{center}
{\Huge Heidelberg in Christian Krachts \textit{„Faserland“}}\\
{\Large Quin Barnard}
\end{center}
%
\begin{flushright}
    Quartal 4 2026\\
    Woche 4\\
    Oktober\\
\end{flushright}
%
\textbf{\large Heidelberg}
\begin{itemize}
    \item Baden-Württemberberg
        \begin{itemize}[label=$\circ$]
            \item Fünfgrößte Stadt im Bundesland
        \end{itemize}
    \item Fluß heißt „der Neckar“ („die Aue“=Floodplain/Vloedvlakte)
    \item Älteste Universität Deutschlands
\end{itemize}
%
\begin{center}
{\large
    \textit{
        „Na, jedenfalls mache ich mit dem Schlüssel die Hoteltür auf, ich bin ja immer noch im Hotel Alt Heidelberg, während ich mich an das blöde Spiel auf Madeira erinnere, und dann gehe ich ins Zimmer und mache erst mal das Licht an, weil es ziemlich düster ist da drinnen.“
    }
}
[S. 82]
\end{center}
%
\vspace{.5em}
\textbf{\large Die Erinnerung als Betäubung}
\begin{itemize}
    \item Die Gegenwart verschwindet völlig
    \item Der Protagonist muss sich aktiv merken, dass er kein Kind mehr ist — Warum?
\end{itemize}
%
\vspace{.5em}
\textbf{\large Die Erinnerung als Sehnen}
\begin{itemize}
    \item \underline{„ich bin volkommen zielstrebig auf den Pool zugelaufen“}
        \begin{itemize}[label=$\circ$]
            \item Das Ziel an sich ist wichtig
            \item Mangelt dem heutigen Protagonist
        \end{itemize}
    \item Das Spiel gab ihm ein Zweck, deswegen will er dahin zurück
\end{itemize}
%
\vspace{.5em}
\textbf{\large Die Erinnerung als Gemeinschaft}
\begin{itemize}
    \item Auch als Kind achtet er nicht auf Andere — „Die Person, an die ich mich überhaupt nicht mehr erinnere“ — doch ist er unmittelbar mit jemandem verbunden
    \item „Markenfetischismus“ = Pseudogemeinschaft: Der Protagonist braucht Mitteln dazu, Sozialverhältnisse zu interpretieren
    \item Ambivalenz = Erkenntis des falschen Eindrucks der Gemeinschaft?
\end{itemize}
%
\clearpage
\medskip\noindent\textbf{\Large Quellen}\\
\onehalfspacing
\setlength{\parindent}{0pt}
{Bühler, Patrick \& Marquardt, Franka: Das »große Nivellier-Land«? Die Schweiz in Christian Krachts \textit{Faserland} In: Birgfeld, Johannes \& Conter, Claude D. (Hgg.): Christian Kracht. \textit{Zu Leben und Werk}. Köln 2009. S.76-91.
\vspace{0.75em}

Hanyin, Zhang: Der Nomade in der Fluchtlinie. Das Werden in Christian Krachts Roman \textit{Faserland}. In: Literaturstraße. Chinesisch-deutsche Zeitschrift für Sprach-und Literaturwissenschaft 22 (2021), H.1.
\vspace{0.75em}

Heidelberg: Deutschland von seiner romantischen Seite. URL:\\ {\url{https://www.germany.travel/de/staedte-kultur/heidelberg.html} [24.09.2026].}
\vspace{0.75em}

Kracht, Christian: Faserland. Köln 1995.
\vspace{0.75em}

Lauer, Julia: Bombenangriffe auf Heidelberg während des Zweiten Weltkriegs. In: Jahhrbuch zur Geschichte der Stadt 26 (2022). URL: \\{\url{https://doi.org/10.57962/regionalia-24428}} [24.09.2026].}

 
\end{document}

